% =========================================================
%  sesi/01/sesi.tex
%  Sesi 1 — Pengenalan Unity dan scene pertama
% =========================================================

\sesi{01}{Pengenalan Unity dan Scene Pertama}{3.1}{100 menit}

Modul ini dimulai sesudah UTS. Pada Blok I dan Blok II Anda membangun
\NamaPermainan\ dengan Python: siklus permainan, gerak, fisika, tabrakan, dan
AI. Sekarang pekerjaan yang sama dipindahkan ke \istilah{game engine} Unity.

Yang dibangun bukan permainan baru. Ini tetap \NamaPermainan, dengan cerita
yang sama: pemain masuk gudang, mengumpulkan artefak, menghindari bahaya,
lalu mencapai pintu keluar. Level yang dikerjakan pada modul ini adalah Level
Rak, yaitu level bertampilan samping.

Sesi pertama belum membuat pemain bergerak. Sesi ini menyiapkan project,
mengenal editor, dan menyusun scene berisi bentuk-bentuk primitif sebagai
pengganti sementara.

% ---------------------------------------------------------
\bagian{Tujuan Sesi}

\begin{tujuanbox}
Setelah menuntaskan sesi ini, Anda dapat menunjukkan:
\begin{enumerate}[leftmargin=1.5em, itemsep=0.2em]
  \item project 2D bernama \NamaProject\ yang terbuka di Unity Editor;
  \item scene \file{LevelRak} berisi kamera, lantai, tiga rak, pemain, dan
        pintu;
  \item susunan \istilah{Hierarchy} yang rapi dengan objek berada di bawah
        induk \code{Gudang};
  \item scene yang tersimpan dan dapat dibuka kembali;
  \item penjelasan padanan konsep Unity dengan konsep Pygame yang telah Anda
        pelajari.
\end{enumerate}
\end{tujuanbox}

% ---------------------------------------------------------
\bagian{Prasyarat}

\begin{prasyaratbox}
\begin{itemize}[leftmargin=1.4em, itemsep=0.25em]
  \item Anda telah menuntaskan Blok I dan Blok II, atau setidaknya memahami
        \istilah{game loop}, \istilah{delta time}, dan \istilah{collision}
        dari Python.
  \item Komputer memenuhi kebutuhan Unity LTS: Windows 10/11 64-bit,
        RAM minimal 8 GB, dan ruang disk sekitar 10 GB.
  \item Unity Hub sudah terpasang. Bila belum, pemasangannya dikerjakan pada
        Langkah 1.
\end{itemize}
\end{prasyaratbox}

% ---------------------------------------------------------
\bagian{Peta Kemampuan}

\petasesi
  {Belum ada project Unity. Kemampuan dari Blok I--II masih dalam bentuk
   kode Python.}
  {Project 2D, scene, kamera, objek primitif, dan skrip C\# pertama.}
  {Level Rak yang dapat dimainkan: \istilah{Tilemap}, \istilah{Collider2D},
   \istilah{Rigidbody2D}, dan kontrol pemain pada sesi 2.}

% ---------------------------------------------------------
\bagian{Konsep yang Dipakai}

\subsection*{Empat panel yang akan sering Anda pakai}

\begin{itemize}[leftmargin=1.4em, itemsep=0.3em]
  \item \textbf{Hierarchy} --- daftar seluruh \istilah{GameObject} pada scene
        yang sedang dibuka. Objek di bawah objek lain menjadi anaknya.
  \item \textbf{Scene} dan \textbf{Game} --- tampilan kerja di editor dan
        tampilan hasil ketika permainan berjalan.
  \item \textbf{Inspector} --- isi objek yang sedang dipilih: transform,
        komponen, dan nilai-nilainya.
  \item \textbf{Project} --- seluruh berkas project: scene, skrip, gambar,
        dan bunyi.
\end{itemize}

\subsection*{GameObject dan komponen}

Di Unity tidak ada kelas yang menggambar dirinya sendiri seperti
\code{pygame.sprite.Sprite}. Sebagai gantinya, sebuah \istilah{GameObject}
adalah wadah kosong; perilaku dan tampilannya datang dari komponen yang
ditempelkan. \code{Sprite Renderer} membuatnya terlihat,
\code{Collider2D} membuatnya dapat ditabrak, dan \code{Rigidbody2D} membuatnya
tunduk pada fisika. Skrip C\# juga komponen.

Karena itu, pertanyaan ``di mana kode gambar untuk pemain?'' berubah menjadi
``komponen apa yang ditempelkan pada GameObject Pemain?''.

\subsection*{Padanan dengan Pygame}

{\small
\renewcommand{\arraystretch}{1.35}
\begin{tabularx}{\linewidth}{@{}L{4.6cm} Y@{}}
\toprule
\textbf{Di Pygame} & \textbf{Di Unity} \\
\midrule
Jendela dari \code{set_mode} & Tampilan \menu{Game} pada scene \\
Objek \code{Sprite} & \istilah{GameObject} + komponen \\
Posisi \code{Vector2} pada \code{rect} & \code{Transform.position} \\
Gambar permukaan & \code{Sprite Renderer} \\
\code{update()} pada siklus sendiri & \code{Update()} yang dipanggil mesin \\
\code{delta_time} dihitung sendiri & \code{Time.deltaTime} \\
Scene buatan sendiri di kode & Berkas \file{.unity} yang disunting di editor \\
\bottomrule
\end{tabularx}
}

Perbedaan pentingnya: siklus permainan, pengukuran waktu, dan penggambaran
tidak lagi Anda tulis. Mesin yang menjalankannya. Tugas Anda berpindah ke
menyusun objek, memasang komponen, dan menulis perilaku.

% ---------------------------------------------------------
\bagian{Tutorial}

Nama menu dan jendela ditulis persis seperti yang tampak pada Unity LTS.
Bila versi Unity Anda sedikit berbeda, cari nama yang paling mirip; jangan
mengubah urutan pekerjaan.

\begin{langkah}
  \item \textbf{Pasang Unity.} Buka Unity Hub, pilih \menu{Installs}, lalu
        \menu{Install Editor}. Pilih versi \textbf{LTS} terbaru dan
        pastikan modul \menu{Windows Build Support} ikut terpasang.

  \item \textbf{Buat project.} Pada Unity Hub pilih \menu{New project},
        pilih templat \menu{2D}, isi nama \NamaProject, tentukan lokasi
        penyimpanan, lalu tekan \menu{Create project}. Tunggu sampai editor
        terbuka.

  \item \textbf{Simpan scene.} Tekan \menu{Ctrl+S}. Simpan sebagai
        \file{Assets/Scenes/LevelRak}. Beri nama folder \file{Scenes} bila
        belum ada. Scene yang belum disimpan akan hilang ketika editor
        ditutup.

  \item \textbf{Atur kamera.} Pilih objek \code{Main Camera} pada
        \menu{Hierarchy}. Pada \menu{Inspector}, setel
        \menu{Projection} menjadi \menu{Orthographic}, \menu{Size} menjadi
        5, dan \menu{Background} menjadi warna \code{\#18203A}. Posisi
        kamera: \code{(0, 0, -10)}.

  \item \textbf{Buat induk Gudang.} Klik kanan pada \menu{Hierarchy}, pilih
        \menu{Create Empty}, beri nama \code{Gudang}. Seluruh objek level
        akan menjadi anaknya agar Hierarchy tetap rapi.

  \item \textbf{Buat lantai.} Klik kanan \code{Gudang}, pilih
        \menu{2D Object > Sprites > Square}, beri nama \code{Lantai}.
        Pada \menu{Inspector}: \code{Transform.position} $(0,-3.5,0)$;
        \code{Sprite Renderer > Draw Mode} menjadi \menu{Tiled};
        \code{Size} $(16,1)$; \code{Color} \code{\#46587E}.

  \item \textbf{Buat tiga rak.} Ulangi langkah 6 tiga kali dengan nama
        \code{RakKiri}, \code{RakTengah}, \code{RakKanan}. Ukuran masing-masing
        $(3,0.4)$, warna \code{\#60749E}. Posisi:
        $(-4,-1.8)$, $(0.5,-0.4)$, $(4.5,-1.8)$.

  \item \textbf{Buat pemain dan pintu.} Tambahkan dua Square lagi:
        \code{Pemain} berukuran $(0.8,0.8)$ berwarna \code{\#2E75B6} pada
        posisi $(-6.5,-3)$; \code{Pintu} berukuran $(1,2)$ berwarna
        \code{\#3CA06E} pada posisi $(7,-2.6)$.

  \item \textbf{Buat folder skrip.} Pada jendela \menu{Project}, klik kanan
        folder \file{Assets}, pilih \menu{Create > Folder}, beri nama
        \code{Scripts}.

  \item \textbf{Buat skrip pertama.} Klik kanan folder \file{Scripts}, pilih
        \menu{Create > C\# Script}, beri nama \code{SkripPertama}. Buka
        berkasnya dan salin isi pada bagian Kode Sesi Ini. Perhatikan bahwa
        nama kelas harus sama dengan nama berkas.

  \item \textbf{Pasang skrip pada objek uji.} Buat Square baru bernama
        \code{CobaSkrip} pada posisi $(0,1.5)$, ukuran $(0.6,0.6)$, warna
        \code{\#DEB242}. Seret berkas \code{SkripPertama} dari jendela
        \menu{Project} ke objek itu, atau tekan \menu{Add Component} lalu
        cari \code{SkripPertama}.

  \item \textbf{Jalankan scene.} Tekan \menu{Play}. Gerakkan objek
        \code{CobaSkrip} dengan tombol panah dan amati jendela
        \menu{Console}. Tekan \menu{Play} lagi untuk berhenti.
\end{langkah}

\begin{ujibox}
Setelah Langkah 12, periksa lima hal berikut.
\begin{itemize}[leftmargin=1.4em, itemsep=0.2em]
  \item Hierarchy memuat \code{Gudang} dengan lima anak: lantai, tiga rak,
        pemain, dan pintu.
  \item Jendela \menu{Game} menampilkan lantai, tiga rak, pemain biru, dan
        pintu hijau.
  \item \code{CobaSkrip} bergerak saat tombol panah ditekan dan berputar
        pelan.
  \item \menu{Console} menampilkan ``Skrip mulai.'' dan laporan setiap dua
        detik.
  \item Setelah Play dihentikan, posisi objek kembali seperti sebelum
        dijalankan.
\end{itemize}
\end{ujibox}

% ---------------------------------------------------------
\bagian{Kode Sesi Ini}

Skrip berikut adalah skrip C\# pertama Anda. Ia belum mengatur pemain; ia
hanya membuktikan bahwa scene berjalan, bahwa \code{Update()} dipanggil
setiap bingkai, dan bahwa \code{Time.deltaTime} tersedia dari mesin.

\tampilkode[language=CSharp]{code/sesi-01/Scripts/SkripPertama.cs}
  {SkripPertama.cs: skrip C\# pertama pada project}

\subsection*{Tabel objek scene}

Tabel berikut adalah daftar objek yang harus ada pada scene Anda. Nilai yang
tidak disebut boleh dibiarkan bawaan.

{\footnotesize
\renewcommand{\arraystretch}{1.3}
\begin{tabularx}{\linewidth}{@{}L{2.3cm} L{3.0cm} Y@{}}
\toprule
\textbf{GameObject} & \textbf{Komponen} & \textbf{Nilai penting} \\
\midrule
Kamera & Transform, Camera &
  Position $(0,0,-10)$; Orthographic; Size 5; Background \code{\#18203A} \\
Lantai & Transform, Sprite Renderer &
  Position $(0,-3.5,0)$; Tiled $(16,1)$; \code{\#46587E} \\
RakKiri & Transform, Sprite Renderer &
  $(-4,-1.8,0)$; $(3,0.4)$; \code{\#60749E} \\
RakTengah & Transform, Sprite Renderer &
  $(0.5,-0.4,0)$; $(3,0.4)$; \code{\#60749E} \\
RakKanan & Transform, Sprite Renderer &
  $(4.5,-1.8,0)$; $(3,0.4)$; \code{\#60749E} \\
Pemain & Transform, Sprite Renderer &
  $(-6.5,-3,0)$; $(0.8,0.8)$; \code{\#2E75B6} \\
Pintu & Transform, Sprite Renderer &
  $(7,-2.6,0)$; $(1,2)$; \code{\#3CA06E} \\
CobaSkrip & Transform, Sprite Renderer, SkripPertama &
  $(0,1.5,0)$; $(0.6,0.6)$; \code{\#DEB242}; Move Speed 3 \\
\bottomrule
\end{tabularx}
}

\checkpoint{01}{Folder checkpoint memuat \file{Scripts/SkripPertama.cs} dan
\file{scene.md}. Scene tidak disertakan sebagai berkas biner; susunlah
mengikuti tabel di atas.}

% ---------------------------------------------------------
\bagian{Kesalahan Umum}

\begin{galatbox}
\begin{itemize}[leftmargin=1.4em, itemsep=0.35em]
  \item \textbf{Objek tidak terlihat di jendela Game.}
        Gejala: objek ada di Hierarchy tetapi layar kosong. Periksa posisi
        $z$ objek (harus 0) dan posisi kamera; pada tampilan 2D, kamera
        berada di $z=-10$ dan melihat ke arah $+z$.
  \item \textbf{Objek terlihat tetapi sangat kecil atau sangat besar.}
        Gejala: ukuran tidak sesuai tabel. Periksa \code{Transform.scale}
        dan \code{Sprite Renderer > Size}; keduanya memengaruhi ukuran
        akhir.
  \item \textbf{Berkas skrip tidak dapat ditempel.}
        Gejala: nama kelas dan nama berkas berbeda, atau berkas tidak berada
        di dalam folder \file{Assets}. Keduanya membuat Unity menolak
        komponen.
  \item \textbf{Perubahan hilang setelah Play dihentikan.}
        Gejala: posisi objek kembali seperti semula. Itu memang perilaku
        Unity: perubahan saat Play tidak disimpan. Ubah scene ketika Play
        berhenti.
  \item \textbf{Console merah saat Play.}
        Gejala: skrip tidak berjalan sama sekali. Baca pesan pertama pada
        Console; biasanya nama kelas salah atau titik koma hilang.
\end{itemize}
\end{galatbox}

% ---------------------------------------------------------
\bagian{Lembar Kerja}

Isi tabel berikut, lalu simpan sebagai bagian dari catatan praktikum Anda.

{\small
\renewcommand{\arraystretch}{1.4}
\begin{tabularx}{\linewidth}{@{}C{0.8cm} L{5.4cm} Y@{}}
\toprule
\textbf{No} & \textbf{Hasil yang dicatat} & \textbf{Keterangan Anda} \\
\midrule
1 & Versi Unity dan templat project yang dipakai & \\[0.55cm]
2 & Jumlah GameObject pada Hierarchy dan susunannya & \\[0.55cm]
3 & Ukuran kamera (orthographic size) dan warna latar & \\[0.55cm]
4 & Isi Console setelah Play selama 10 detik & \\[0.9cm]
5 & Satu kesalahan yang muncul dan cara Anda memperbaikinya & \\[0.9cm]
\bottomrule
\end{tabularx}
}

% ---------------------------------------------------------
\bagian{Checklist Sesi}

\begin{checklistbox}
\begin{checklist}
  \item Saya dapat menunjukkan project 2D \NamaProject\ yang terbuka dari
        Unity Hub.
  \item Saya dapat menjelaskan fungsi Hierarchy, Inspector, Project, Scene,
        dan Game.
  \item Saya dapat membuat GameObject beserta komponen Sprite Renderer dan
        mengatur warnanya.
  \item Saya dapat menyimpan scene dan membukanya kembali.
  \item Saya dapat menjelaskan padanan \istilah{GameObject} dengan
        \code{Sprite} pada Pygame.
  \item Saya dapat membaca Console dan menemukan penyebab skrip tidak
        berjalan.
\end{checklist}
\end{checklistbox}

% ---------------------------------------------------------
\bagian{Pengait ke Sesi Berikutnya}

Scene ini masih berupa bentuk primitif yang tidak bertabrakan dan tidak dapat
dikendalikan. Pada sesi 2 Anda menyusun lantai dan rak dengan
\istilah{Tilemap}, menambahkan \istilah{Collider2D} dan
\istilah{Rigidbody2D}, lalu menulis skrip kontrol pemain. Di sanalah
\NamaPermainan\ mulai dapat dimainkan.
